% RQ4 Discussion: Treasury Resilience

\subsection{Interpretation of Results}

\subsubsection{The 10-20\% Buffer Sweet Spot}

Our simulations identify a practical sweet spot for buffer fractions: 10-20\% provides substantial volatility reduction without excessive capital lockup. This aligns with corporate finance norms where firms typically maintain 1-3 months of operating expenses as cash reserves.

Below 10\%, treasuries remain vulnerable to spending shocks and market downturns. Above 20\%, additional buffer provides diminishing protection while increasingly constraining operational flexibility.

The optimal buffer depends on:
\begin{itemize}
    \item Revenue predictability (more volatile revenue $\rightarrow$ higher buffer)
    \item Spending flexibility (can spending be quickly reduced if needed?)
    \item Market regime expectations (bear market $\rightarrow$ higher buffer)
    \item Risk tolerance (conservative organizations prefer higher buffers)
\end{itemize}

\subsubsection{Spending Limits: Sustainability vs. Agility}

Spending limits present a clear tradeoff. Aggressive limits (2\% per period) prevent treasury depletion but may:
\begin{itemize}
    \item Block time-sensitive opportunities
    \item Create proposal backlogs
    \item Fragment initiatives across multiple funding cycles
\end{itemize}

Moderate limits (5\%) balance sustainability with operational flexibility. No limits invite treasury depletion during periods of high proposal activity.

We recommend:
\begin{itemize}
    \item Base limits calibrated to sustainable burn rate
    \item Override mechanisms for urgent/strategic initiatives
    \item Periodic review of limit adequacy
\end{itemize}

\subsubsection{Emergency Mechanisms as Insurance}

Top-up mechanisms function as insurance: they impose cost (reduced spending, diverted revenue) but protect against tail risk (severe drawdown or depletion).

Revenue diversion is gentler---it preserves operational continuity while gradually restoring buffer. Spending freeze is more aggressive---it rapidly restores buffer but disrupts operations.

The choice depends on:
\begin{itemize}
    \item Revenue stability (stable revenue enables gradual diversion)
    \item Spending criticality (discretionary vs. essential spending)
    \item Recovery urgency (how quickly must buffer be restored?)
\end{itemize}

\subsection{Policy Design Framework}

Based on our results, we propose a treasury policy framework:

\subsubsection{Tier 1: Baseline Policies}

\begin{itemize}
    \item Buffer fraction: 15-20\% of treasury
    \item Spending limit: 5\% per period (monthly or quarterly)
    \item Top-up trigger: When buffer falls below 50\% of target
\end{itemize}

\subsubsection{Tier 2: Market-Adaptive}

\begin{itemize}
    \item Dynamic buffer: +5\% in bear markets, -5\% in bull markets
    \item Counter-cyclical spending: Tighter limits during drawdowns
    \item Regime monitoring: Track market conditions for policy triggers
\end{itemize}

\subsubsection{Tier 3: Strategic Reserves}

\begin{itemize}
    \item Separate strategic reserve for multi-year initiatives
    \item Diversification requirements (stablecoin minimums)
    \item External backstops (protocol fee claims, insurance)
\end{itemize}

\subsection{Comparison with Real Practices}

\subsubsection{MakerDAO}

MakerDAO's stability buffer and surplus mechanism parallel our top-up approach. Our simulations validate the value of such mechanisms while suggesting explicit buffer targets.

\subsubsection{Optimism}

Optimism's structured allocations (ecosystem fund, governance fund) implement implicit spending limits by purpose. Our results support this approach while suggesting explicit per-period caps.

\subsubsection{Most DAOs}

Most DAOs lack formal treasury policies. Our results demonstrate the cost of this oversight: increased volatility, longer drawdowns, and higher depletion risk.

\subsection{Limitations}

\subsubsection{Simplified Market Model}

Our market simulation uses stylized regimes. Real crypto markets exhibit:
\begin{itemize}
    \item Fat tails (extreme events more common than modeled)
    \item Correlation spikes (assets correlate during crises)
    \item Liquidity constraints (large treasury sales affect prices)
\end{itemize}

\subsubsection{Exogenous Revenue}

We model revenue as exogenous. In reality, treasury policies affect protocol health, which affects revenue---a feedback loop not fully captured.

\subsubsection{No Strategic Behavior}

We do not model strategic proposer behavior in response to treasury policies (e.g., front-running spending limits, timing proposals to market conditions).

\subsection{Future Work}

\begin{enumerate}
    \item \textbf{Diversification analysis}: Model multi-asset treasury composition and rebalancing
    \item \textbf{Protocol revenue modeling}: Integrate treasury health with protocol usage dynamics
    \item \textbf{Cross-DAO comparison}: Empirical validation against real treasury data
    \item \textbf{Insurance products}: Model external insurance and hedging strategies
\end{enumerate}
